% The Embedding into CP^2  (source §8, §9; added: holomorphy of Phi at O, the
% flex at infinity, the isomorphism-of-Riemann-surfaces theorem proved from
% injectivity + compactness, and inverses as reflection)
\chapter{The Embedding into \texorpdfstring{$\CP^2$}{CP2}}
\label{ch:embedding}

With \(\wp\) and \(\wp'\) in hand we can build the map of \cref{thm:main}: the
embedding
\[
  \Phi : T = \C/\Lambda \longrightarrow E \subset \CP^2,
  \qquad \Phi([z]) = [\wp(z) : \wp'(z) : 1] .
\]
This chapter shows \(\Phi\) is a well-defined holomorphic map onto the smooth
cubic \(E\), extends it across the origin to the point at infinity \(O =
[0:1:0]\), and proves it is an isomorphism of compact Riemann surfaces --- the
analytic half of \cref{thm:main}. Finally we read off the first piece of the
group dictionary: the torus inverse \(z \mapsto -z\) becomes reflection across
the \(x\)-axis. The group law itself is transported in \cref{ch:transport}.

\section{The Projective Map}
\label{sec:emb-map}

For \(z \notin \Lambda\) set \(\Phi([z]) = [\wp(z) : \wp'(z) : 1]\). Since \(\wp,
\wp'\) are \(\Lambda\)-periodic (\cref{prop:wp-even-periodic}) this is
well-defined on \(T\). The differential equation places the image on the cubic.

\begin{proposition}
\label{prop:emb-lands}
For \(z \notin \Lambda\), the point \(\Phi([z])\) lies on
\[
  E : \quad Y^2 Z = 4 X^3 - g_2 X Z^2 - g_3 Z^3 .
\]
\end{proposition}

\begin{proof}
In the affine chart \(Z = 1\) the equation reads \(y^2 = 4x^3 - g_2 x - g_3\).
Substituting \(x = \wp(z)\), \(y = \wp'(z)\) gives exactly the Weierstrass
differential equation \(\wp'^2 = 4\wp^3 - g_2\wp - g_3\) of \cref{thm:wp-ode},
which holds; so \((\wp(z), \wp'(z)) \in E\).
\end{proof}

\begin{figure}[ht]
\centering
\begin{tikzpicture}[scale=1.0, every node/.style={font=\small}]
  \begin{scope}[scale=0.65]
    \drawprecisetorus
    \fill[formteal!85!black] (\torusMarkZx,\torusMarkZy) circle (2.5pt);
    \node[formteal!60!black, right, font=\scriptsize] at (\torusMarkZx,\torusMarkZy) {$[z]$};
    \fill[black] (\torusMarkOriginx,\torusMarkOriginy) circle (2.5pt);
    \node[font=\scriptsize, above] at (\torusMarkOriginx,\torusMarkOriginy) {$[0]$};
    \node at (0,-1.9) {\textbf{torus }$T=\C/\Lambda$};
  \end{scope}
  \draw[-{Latex[length=3mm]}, formplum, line width=1.1pt] (2.0,0) -- (3.4,0);
  \node[formplum] at (2.7,0.4) {$\Phi$};
  \node[formplum, font=\scriptsize, align=center] at (2.7,-0.4) {$[z]\mapsto$\\$[\wp(z){:}\wp'(z){:}1]$};
  \begin{scope}[xshift=6.0cm]
    \draw[->, black!50, line width=0.4pt] (-1.4,0) -- (2.1,0) node[right, black!55, font=\scriptsize] {$x$};
    \draw[->, black!50, line width=0.4pt] (0,-1.6) -- (0,1.6) node[above, black!55, font=\scriptsize] {$y$};
    \draw[draw=formteal!60!black, line width=0.85pt, smooth, samples=70, variable=\x, domain=-1.0:0.0]
      plot (\x, {0.65*sqrt(max(0,4*\x*\x*\x-4*\x))});
    \draw[draw=formteal!60!black, line width=0.85pt, smooth, samples=70, variable=\x, domain=-1.0:0.0]
      plot (\x, {-0.65*sqrt(max(0,4*\x*\x*\x-4*\x))});
    \draw[draw=formteal!60!black, line width=0.85pt, smooth, samples=70, variable=\x, domain=1.0:1.9]
      plot (\x, {0.65*sqrt(max(0,4*\x*\x*\x-4*\x))});
    \draw[draw=formteal!60!black, line width=0.85pt, smooth, samples=70, variable=\x, domain=1.0:1.9]
      plot (\x, {-0.65*sqrt(max(0,4*\x*\x*\x-4*\x))});
    \foreach \rt in {\rooteone,\rootetwo,\rootethree}
      \fill[formteal!60!black] (\rt,0) circle (1.2pt);
    \fill[formteal!85!black] (\chordQx,{0.65*\chordQy}) circle (1.8pt);
    \node[formteal!60!black, right, font=\scriptsize] at (\chordQx,{0.65*\chordQy}) {$\Phi([z])$};
    \node at (0.4,-1.85) {\textbf{cubic} $E\subset\CP^2$};
    \node[font=\scriptsize] at (0.4,1.85) {$y^2=4x^3-4x$};
  \end{scope}
\end{tikzpicture}
\caption{The map \(\Phi([z]) = [\wp(z):\wp'(z):1]\) carries the torus onto the
real locus of the cubic \(y^2 = 4x^3 - 4x\) (here \(g_2 = 4\), \(g_3 = 0\)).
Curve and sample point are from \texttt{sagemath/12\_curve\_sample.sage}.}
\label{fig:embedding}
\end{figure}

\section{Extension at the Origin}
\label{sec:emb-origin}

At \(z = 0\) both \(\wp\) and \(\wp'\) blow up, so the affine formula fails, but
the projective one has a limit. Multiplying all three homogeneous coordinates by
\(z^3\),
\[
  [\wp(z) : \wp'(z) : 1] = [z^3\wp(z) : z^3\wp'(z) : z^3],
\]
and using \(\wp = z^{-2} + O(z^2)\), \(\wp' = -2z^{-3} + O(z)\),
\[
  z^3\wp(z) \to 0, \qquad z^3\wp'(z) \to -2, \qquad z^3 \to 0
  \quad (z \to 0).
\]
Hence \(\lim_{z\to0}\Phi([z]) = [0 : -2 : 0] = [0:1:0]\), and we set
\[
  \Phi(O) = [0:1:0] =: O \in E .
\]

\begin{proposition}
\label{prop:emb-holo}
With this value at the origin, \(\Phi : T \to E\) is holomorphic.
\end{proposition}

\begin{proof}
Away from \([0]\), \(\Phi\) is given by the holomorphic functions \(\wp, \wp'\) in
the chart \(Z = 1\), so it is holomorphic there. Near \([0]\) use the homogeneous
representative \(\bigl(z^3\wp(z),\, z^3\wp'(z),\, z^3\bigr)\): each entry is
holomorphic in \(z\) (the poles of \(\wp, \wp'\) are cancelled by \(z^3\)), and
they do not vanish simultaneously, since the middle entry tends to \(-2 \ne 0\).
In the affine chart \(Y = 1\) of \(\CP^2\) the map is therefore
\(\bigl(z^3\wp/(z^3\wp'),\ z^3/(z^3\wp')\bigr) = \bigl(\wp/\wp',\ 1/\wp'\bigr)\),
holomorphic at \(z = 0\) with value \((0,0) = O\). So \(\Phi\) is holomorphic on
all of \(T\).
\end{proof}

\section{The Point at Infinity}
\label{sec:emb-infinity}

\begin{proposition}
\label{prop:emb-flex}
\(O = [0:1:0]\) is the unique point of \(E\) on the line at infinity \(Z = 0\). It
is a smooth point of \(E\), and the line at infinity is the tangent there, meeting
\(E\) only at \(O\) with intersection multiplicity three: \(O\) is an inflection
point.
\end{proposition}

\begin{proof}
Setting \(Z = 0\) in \(Y^2 Z = 4X^3 - g_2 X Z^2 - g_3 Z^3\) gives \(0 = 4X^3\), so
\(X = 0\); the only such projective point is \([0:1:0]\). In the affine chart
\(Y = 1\), with \(F(X,Z) = Z - 4X^3 + g_2 X Z^2 + g_3 Z^3\), we have
\(\partial F/\partial Z\,|_{(0,0)} = 1 \ne 0\), so \(O\) is smooth and the tangent
is \(Z = 0\). Substituting \(Z = 0\) into \(E\) yields \(4X^3 = 0\), a triple root
at \(X = 0\), so \(Z = 0\) meets \(E\) at \(O\) with multiplicity three.
\end{proof}

The flex at \(O\) is the geometric reason \(O\) will serve as the identity of the
group law: a line meeting \(E\) three times at \(O\) records \(3O\), the divisor
class that is the identity (\cref{ch:transport}).

\section{Why \texorpdfstring{$L(3O)$}{L(3O)} Gives the Embedding}
\label{sec:emb-l3o}

The coordinates of \(\Phi\) are a basis of \(L(3O)\). By
\cref{prop:rr-counts}, \(\ell(3O) = 3\), with basis \(\{1, \wp, \wp'\}\): a function
with no pole, one with a double pole, one with a triple pole, all at \(O\). Three
linearly independent functions define a map to \(\CP^2\),
\[
  [z] \longmapsto [\wp(z) : \wp'(z) : 1],
\]
and the single algebraic relation among the three sections --- forced by the
dimension count of \cref{rem:rr-cubic-preview} --- is the cubic. Because
\(\Delta \ne 0\) (\cref{cor:wp-smooth}), \(E\) is a \emph{smooth} projective plane
curve, hence a compact Riemann surface: near each point one of the partial
derivatives of the defining equation is nonzero, and the implicit function
theorem provides a holomorphic chart. It is also connected (the polynomial
\(Y^2 - (4X^3 - g_2 X - g_3)\) is irreducible when the cubic has distinct roots).

\section{Injectivity and the Isomorphism}
\label{sec:emb-iso}

\begin{lemma}
\label{lem:emb-fibers}
If \(\wp(z) = \wp(w)\) then \(w \equiv z\) or \(w \equiv -z \pmod\Lambda\).
\end{lemma}

\begin{proof}
Fix \(w\) and consider \(\wp - \wp(w)\), elliptic of order two (double pole at
\([0]\)), so by \cref{thm:ell-valence} it has exactly two zeros counted with
multiplicity. Both \(w\) and \(-w\) are zeros, since \(\wp\) is even; if
\(w \not\equiv -w\) these are the two distinct zeros, and \(\wp(z) = \wp(w)\)
forces \(z \equiv w\) or \(z \equiv -w\). If \(w \equiv -w\) (a half-period), then
\(\wp - \wp(w)\) has a double zero at \(w\), again accounting for both zeros, and
\(\wp(z) = \wp(w)\) forces \(z \equiv w \equiv -w\). In all cases
\(z \equiv \pm w\).
\end{proof}

\begin{theorem}[\texorpdfstring{$\Phi$}{Phi} is an isomorphism]
\label{thm:emb-iso}
The map \(\Phi : T \to E\) is a biholomorphism of compact Riemann surfaces.
\end{theorem}

\begin{proof}
\emph{Injectivity.} Suppose \(\Phi([z]) = \Phi([w])\) with neither at \(O\). Then
\(\wp(z) = \wp(w)\) and \(\wp'(z) = \wp'(w)\). By \cref{lem:emb-fibers},
\(w \equiv z\) or \(w \equiv -z\). In the second case \(\wp'(w) = \wp'(-z) =
-\wp'(z)\), and combined with \(\wp'(w) = \wp'(z)\) this gives \(\wp'(z) = 0\),
i.e.\ \(z\) is a half-period (\cref{prop:wp-roots}), where \(z \equiv -z\); so
\(w \equiv z\) anyway. As \(O\) is the only point mapping to \([0:1:0]\)
(\cref{prop:emb-flex}), \(\Phi\) is injective.

\emph{Bijectivity and biholomorphy.} \(\Phi\) is a nonconstant holomorphic map
between the compact connected Riemann surfaces \(T\) and \(E\); its image is open
(open mapping theorem) and compact, hence all of the connected curve \(E\), so
\(\Phi\) is surjective. An injective holomorphic map between Riemann surfaces is
an immersion with nonvanishing derivative, so \(\Phi\) is a holomorphic bijection
whose inverse is holomorphic. Therefore \(\Phi\) is a biholomorphism.
\end{proof}

This proves the Riemann-surface assertion of \cref{thm:main}: the torus and the
cubic are the same compact Riemann surface, identified by \((\wp, \wp')\). The
remaining assertion --- that \(\Phi\) is a group isomorphism for the chord--tangent
law --- is the subject of \crefrange{ch:abel}{ch:transport}.

\section{Inverses Become Reflection}
\label{sec:emb-reflection}

Transport the torus inverse \(z \mapsto -z\) through \(\Phi\). Since \(\wp\) is
even and \(\wp'\) odd,
\[
  \Phi(-z) = [\wp(-z) : \wp'(-z) : 1] = [\wp(z) : -\wp'(z) : 1],
\]
so in affine coordinates \(P = (x, y) = (\wp(z), \wp'(z))\) maps to
\[
  -P = (x, -y) .
\]

\begin{figure}[ht]
\centering
\begin{tikzpicture}[scale=1.0, every node/.style={font=\small}]
  \draw[->, black!50, line width=0.5pt] (-1.8,0) -- (3.0,0) node[right, black!55] {$x=\wp$};
  \draw[->, black!50, line width=0.5pt] (0,-3.4) -- (0,3.5) node[above, black!55] {$y=\wp'$};
  \drawprecisecurve
  \foreach \rt in {\rooteone,\rootetwo,\rootethree}
    \fill[formteal!60!black] (\rt,0) circle (1.3pt);
  \fill[formrust!85!black] (\chordQx,\chordQy) circle (2pt);
  \node[formrust!60!black, right] at (\chordQx,\chordQy) {$P=(\wp(z),\wp'(z))$};
  \fill[formgold!85!black] (\chordQx,{-\chordQy}) circle (2pt);
  \node[formgold!50!black, right] at (\chordQx,{-\chordQy}) {$-P=(\wp(-z),-\wp'(z))$};
  \draw[<->, draw=formplum, line width=0.6pt, dashed]
    (\chordQx,{\chordQy-0.18}) -- (\chordQx,{-\chordQy+0.18});
  \node[formplum, font=\scriptsize, anchor=west] at ({\chordQx+0.10},0) {reflect};
  \node[draw=black!40, rounded corners=2pt, fill=formivory, align=left, anchor=south,
        inner sep=4pt, font=\scriptsize] at (0.0,-3.4)
    {$\wp$ even, $\wp'$ odd $\Rightarrow$ $z\mapsto-z$ on the torus becomes
     $(x,y)\mapsto(x,-y)$ on the curve.\\
     Sage: $P=(\chordQx,\chordQy)$, $-P=(\chordQx,{-\chordQy})$.};
\end{tikzpicture}
\caption{The torus inverse \(z \mapsto -z\) becomes reflection across the
\(x\)-axis on the cubic. The fixed points \(y = 0\) are the half-periods
\(e_1, e_2, e_3\); \(O\) is its own inverse at infinity. Curve from
\texttt{sagemath/12\_curve\_sample.sage}.}
\label{fig:reflection}
\end{figure}

\begin{keyidea}{Why reflection appears}
Reflection across the \(x\)-axis is not an extra convention bolted onto the
chord--tangent rule. It \emph{is} the torus inverse \(z \mapsto -z\), written in
the coordinates \(x = \wp(z)\), \(y = \wp'(z)\). The points fixed by the
reflection are exactly those with \(y = 0\) --- the images
\((e_1, 0), (e_2, 0), (e_3, 0)\) of the half-periods --- together with \(O\) at
infinity. These four are the two-torsion of the group: the points equal to their
own inverse.
\end{keyidea}

\section{Summary}
\label{sec:emb-summary}

The map \(\Phi([z]) = [\wp(z):\wp'(z):1]\), extended by \(\Phi(O) = [0:1:0]\), is a
holomorphic map of \(T\) onto the smooth cubic \(E\)
(\cref{prop:emb-lands,prop:emb-holo}). The point \(O\) is the unique point at
infinity and an inflection (\cref{prop:emb-flex}), destined to be the identity.
The coordinates are a basis of \(L(3O)\), and because \(\Delta \ne 0\) the curve
\(E\) is a compact Riemann surface; \(\Phi\) is injective
(\cref{lem:emb-fibers}) and hence a biholomorphism
\(T \xrightarrow{\sim} E\) (\cref{thm:emb-iso}), the analytic content of
\cref{thm:main}. Under \(\Phi\) the torus inverse \(z \mapsto -z\) is reflection
\((x,y) \mapsto (x,-y)\), whose fixed points are the two-torsion. What remains is
to see that the \emph{addition} on \(T\) becomes the chord--tangent law; that
needs Abel's sum law, taken up next.

\section*{Exercises}
\addcontentsline{toc}{section}{Exercises}

\begin{exercise}
\label{ex:emb-holo-origin}
Verify \cref{prop:emb-holo} in the chart \(Y = 1\): show the map near \([0]\) is
\(z \mapsto (\wp/\wp', 1/\wp')\), compute its value and derivative at \(z = 0\),
and confirm it is a local holomorphic coordinate. (This shows \(\Phi\) is an
immersion at \(O\).)
\end{exercise}

\begin{exercise}
\label{ex:emb-flex}
Carry out \cref{prop:emb-flex} in full: confirm \(O\) is smooth by computing the
gradient of the projective equation there, and verify the intersection
multiplicity of \(Z = 0\) with \(E\) at \(O\) is three by substitution. Why does a
flex at \(O\) make \(3O\) a principal divisor?
\end{exercise}

\begin{exercise}
\label{ex:emb-fiber-lemma}
Prove \cref{lem:emb-fibers} and deduce that \(z \mapsto \wp(z)\) realizes
\(T/\{\pm1\} \cong \CP^1\), with the four fixed points of \(z \mapsto -z\) mapping
to the branch values \(e_1, e_2, e_3, \infty\).
\end{exercise}

\begin{exercise}
\label{ex:emb-two-torsion-image}
Show that \(\Phi\) sends the four two-torsion points to the points of \(E\) fixed
by \((x,y) \mapsto (x,-y)\), according to
\[
  0 \mapsto O, \quad
  \tfrac{\omega_1}{2} \mapsto (e_1,0), \quad
  \tfrac{\omega_2}{2} \mapsto (e_2,0), \quad
  \tfrac{\omega_1+\omega_2}{2} \mapsto (e_3,0),
\]
and that these four points are precisely those of \(E\) with \(y = 0\) or
\(y = \infty\).
\end{exercise}

\begin{exercise}
\label{ex:emb-involution}
Show the reflection \(\iota_E(x,y) = (x,-y)\) is a holomorphic involution of \(E\)
with quotient \(\CP^1\) (via \(x = \wp\)), and that under \(\Phi\) it corresponds
to \(z \mapsto -z\). Conclude that \(\iota_E\) is the unique nontrivial deck
transformation of the double cover \(\wp : T \to \CP^1\).
\end{exercise}
